\begin{corollary}[C.3.27 (Proved): a second determination by the spin$^c$ coboundary]
The actual positive-loop smooth monodromy satisfies
$[\Phi]=0\in\Theta_7$. This determination uses the explicit
eight-dimensional coboundary of C.3.25 and the direct index
additivity of C.3.26. It does not use the three-knot integer,
Wall's formula for it, the Davis generator comparison, or the
conclusion of C.3.24.
\end{corollary}
\begin{proof}
Apply C.3.26 to equivariant representatives of the actual return
maps. For the annular six-loop representative $\psi$, the two
geometric characteristic calculations, with their common
orientation sign, give
\[
 p_\psi=-24\epsilon_{\rm cap},\qquad
 k_\psi=-6\epsilon_{\rm cap}.
\]
These are the actual Pontryagin and cubic pairings proved in
C.3.18--C.3.22. Therefore the directly defined index is
\[
 \lambda(\psi)=\frac{p_\psi-4k_\psi}{24}=0
 \quad\hbox{in }\Z.
\]

We check the relation to one positive loop before reducing
modulo $28$. Choose an invariant connection on the descended
$w$-family. Its return over six consecutive positive loops is
the sixth power of its one-loop return. A positive loop changes
$c$ by $-2$. The coherent closing translations of C.3.5 obey
\[
 b_{h,c}+b_{k,c+h}=b_{h+k,c};
\]
in particular the six successive positive-$2$ closings from
$c_0-12$ to $c_0$ equal the positive-$12$ closing. Interpolating
between invariant connections gives equivariant isotopies.
The changes of transport in C.3.10, its unit and real-scalar
isotopies, and the lift comparisons of C.3.6 are equivariant
as well. Changing back to the stipulated closing choices
adds only a connected $\C^*$ automorphism commuting with the
circle action. Consequently the actual lift of $\psi$ and
$\Phi^6$ define the same equivariant mapping class.

If necessary, a lift of the orbit identity on the fixed tube
is made the identity there by its smooth circle gauge; the
gauge has a real logarithm and can be extended with a collar
cutoff. This is an equivariant isotopy and does not change the
orbit map. C.3.26 therefore applies with exactly the cap used
in the calculation above, and its normalization independence
gives
\[
 6\lambda(\Phi)=\lambda(\Phi^6)=\lambda(\psi)=0.
\]
The group here is $\Z$, so $\lambda(\Phi)=0$.

Now construct the actual eight-dimensional spin$^c$ coboundary
$Z_\Phi$ of C.3.25 for a normalized one-loop representative.
Its explicit intersection form and characteristic covectors
give, by the spin$^c$ index formula alone,
\[
 \mu(\Sigma_\Phi)
   \equiv\frac{p_\Phi-4k_\Phi}{672}
   =\frac{\lambda(\Phi)}{28}=0\pmod\Z.
\]
The positive generator in C.1.2 has invariant $1/28$, and
$\Theta_7$ has order $28$. Thus this invariant detects its
elements and $[\Sigma_\Phi]=0$. The sphere-gluing comparison
of C.1.1 gives $[\Phi]=0$. It also gives
$M_\Phi\cong S^1\times S^6$ with either spin structure having
zero invariant.
\end{proof}

\begin{remark}[Independence and shared geometric inputs]
The first determination uses the integer of equivariant
orbit knots and its Wall--Davis comparison. The present
determination uses the explicit line-disc bundle and
four-handle coboundary, the spin$^c$ index formula, and an
independent integer-additivity argument from actual
seven-dimensional characteristic bordisms. Neither conversion
nor the integer-level passage from six loops to one in this
second argument invokes the first determination. The two
methods share the proved annular representative and the actual
geometric values $p_\psi=-24\epsilon_{\rm cap}$ and
$k_\psi=-6\epsilon_{\rm cap}$. Their independence is at the
boundary-invariant and one-loop arguments; they are not
calculations from disjoint geometric data. Independent
referees and exact arithmetic checks are validation, not
additional mathematical methods.

The smooth spin$^c$ coboundary here is different from the
specific analytic filling proposed in Part B. Its construction
does not supply that filling or its almost-complex Chern data.
\end{remark}
